% Generated from chapter-05.md - do not edit.

\chapter{Thinking About The Three Principles}%
\index{three principles}\nobreak%

If we are going to ask ourselves, ``What's a better thought?'' it helps to understand just how important thought may be in creating our experience of life. The Three Principles offer one model for looking at this---one that places Thought at the very center of how we experience ourselves and the world around us.

\section{A Kinder Way of Understanding How We Experience Life}%
\index{kinder}\nobreak%

Every now and then, someone explains something in a way that feels both simple and profound at the same time. That's how I felt when I first heard about what are called the \textbf{Three Principles} --- originally named \textbf{Mind, Consciousness, and Thought}.

They were introduced by a quiet, thoughtful man named \textbf{Sydney Banks}, who had no formal background in psychology or academia. His realization was simple and life-changing:\\
\textbf{everything we experience in life comes from the inside out, not from the world itself.}\\
What's happening ``out there'' doesn't directly create our feelings --- our inner process does.

That insight opened up a whole new way of seeing life for me.

Traditionally, the Three Principles are described like this:

\begin{itemize}
\item \textbf{Mind} --- the universal intelligence behind life
\item \textbf{Consciousness} --- our capacity to be aware
\item \textbf{Thought} --- the creative power that forms our moment-to-moment experience
\end{itemize}

I honor that language and where it came from.

And---for my own understanding, and for the sake of being grounded and practical---I think about it this way:

\textbf{The Brain, Consciousness, and Thought are working together to create our entire human experience.}

\textbf{BRAIN}

\textbf{For me, the brain is where the miracle of human experience shows up in everyday life.}

I am a living human being with a remarkable brain---an organ so powerful and compassionate that it quietly runs my breathing, digestion, heartbeat, immune system, and nervous system . . . 

and creates my thinking, my emotions, my perceptions, my joys, and my worries.

\emph{(All without asking me to fill out any forms.)}

\section{Consciousness --- The Gift of Awareness}%
\index{consciousness}\index{awareness}\nobreak%

Consciousness is what allows us to be aware at all.

It's the light that brings our thoughts to life.

It's what lets us notice a beautiful sunset, feel comfort from a warm cup of tea, experience fear, or suddenly see something from a new angle.

It's the difference between simply existing . . .  and experiencing.

You know how a dream feels completely real while you're inside it?

That's consciousness at work.

\emph{(Even the very strange dreams.)}

\section{Pause for a Moment}

Right now, what are you aware of?

A sound?\\
A feeling in your body?\\
A thought passing through?

Notice how awareness brings it to life.

And gently consider:

\textbf{``I'm feeling this way because I'm conscious of my thinking---not because life itself is actually terrible.''}

\section{Thought --- The Creative Power of Experience}

Thought is the moment-to-moment creative force of our lives.

We use it constantly---usually without noticing---to create our inner world.

Every feeling, every story, every judgment, every worry, every hope . . . 

we experience them through the lens of thought.

We don't live in the feeling of our circumstances.

We live in the feeling of our thinking about our circumstances.

\emph{(Which explains why two people can experience the same event . . .  and feel completely different.)}

\textbf{This doesn't mean thoughts are bad or that we should try to control them.}

\textbf{It simply means they are always changing---like weather in the sky.}

\emph{(Sometimes sunny, sometimes stormy, occasionally a bit unpredictable.)}

\textbf{Better thought:}

\emph{``This is just a passing thought---not a permanent truth.''}

\section{Why This Is So Comforting}

What comforts me most in understanding all of this is something very gentle:

\textbf{We are not broken.}\\
\textbf{We don't need fixing.}

\textbf{We are simply living inside the ever-changing flow of brain, consciousness, and thought.}

\textbf{Sometimes we create storms in our own inner weather.}

\textbf{And sometimes we forget that the sky always clears.}

\emph{(Even when it takes a little while.)}

\textbf{When we realize that our feelings are coming from the thoughts we're having---not directly from the world---it becomes easier to be kind to ourselves.}

\textbf{To pause.}\\
\textbf{To breathe.}\\
\textbf{To wait for the storm to pass.}

\textbf{And often, all it takes is one small shift---}

\textbf{just one better thought---}

\textbf{to change everything.}

\section{A Gentle Noticing About Beliefs}

I've also begun to notice something important:

\textbf{Thoughts that are repeated often---and felt strongly---can become beliefs.}

\textbf{They start as passing ideas . . . }

\textbf{but over time, they can feel solid, fixed, and true.}

\textbf{And yet---even beliefs began as thoughts.}

\emph{(Very convincing thoughts . . .  but still thoughts.)}

\textbf{Which means:}

\textbf{They can be questioned.}\\
\textbf{They can soften.}\\
\textbf{They can change.}

\section{A Quiet Ending Thought}

You don't have to understand all of this perfectly.

You don't have to ``get it right.''

Just begin to notice:

\begin{itemize}
\item What am I thinking right now?
\item How is that thought shaping how I feel?
\end{itemize}

And if you'd like, gently ask:

\textbf{What's a better thought?}

\emph{Even a small one can begin to shift your whole experience.}
